% Part: computability
% Chapter: recursive-functions
% Section: other-recursions

\documentclass[../../../include/open-logic-section]{subfiles}

\begin{document}

\olfileid{cmp}{rec}{ore}
\olsection{بازگشت دے ہور روپ}

جوڑیاں تے سلسلے بناون دے طریقے ورت کے، بنیادی
بازگشت دے ہور وکھرے (تے مفید) روپاں دا جواز دتا جا
سکدا اے۔ مثال لئی، اکثر دو فنکشناں نوں اکّو ویلے
متعین کرنا مفید ہوندا اے، جداں ہیٹھلی تعریف وچ:
\begin{align*}
h_0(\vec x, 0) & = f_0(\vec x) \\
h_1(\vec x, 0) & = f_1(\vec x) \\
h_0(\vec x, y+1) & = g_0(\vec x, y, h_0(\vec x, y), h_1(\vec x, y)) \\
h_1(\vec x, y+1) & = g_1(\vec x, y, h_0(\vec x, y), h_1(\vec x, y))
\end{align*}
ایہہ \emph{اکٹھّی بازگشت (simultaneous recursion)}
دی اک مثال اے۔ فنکشن متعین کرن دا اک ہور مفید طریقہ
ایہہ اے کہ $h(\vec x, y+1)$ دی قیمت نوں
$h(\vec x, 0)$، \dots،~$h(\vec x, y)$ دیاں
\emph{ساریاں} قیمتاں دے حوالے نال دتا جاوے، جداں
ہیٹھلی تعریف وچ:
\begin{align*}
h(\vec x, 0) & = f(\vec x) \\
h(\vec x, y+1) & = g(\vec x, y, \tuple{h(\vec x, 0), \dots, h(\vec x, y)}).
\end{align*}
ہیٹھلا خاکہ ایس خیال نوں ہور مختصر طور اُتے ظاہر
کردا اے:
\[
h(\vec x, y) = g(\vec x, y, \tuple{h(\vec x, 0), \dots, h(\vec x, y-1)})
\]
ایتھے $g$ دی آخری دلیل بارے ایہہ سمجھنا اے کہ
جدوں $y$، $0$ ہووے، تاں اوہ صرف خالی سلسلہ اے۔ دوہاں
روپاں وچ خیال ایہہ اے کہ «جانشین والا قدم» حساب
کردے ویلے فنکشن $h$ ہن تک حساب کیتیاں قیمتاں
دا پورا سلسلہ ورت سکدا اے۔ ایس نوں
\emph{پچھلی ساریاں قیمتاں والی} بازگشت
(course-of-values recursion) کہندے نیں۔ اک خاص
مثال لئی، ایس نال ہیٹھلی قسم دی تعریف دا جواز دتا
جا سکدا اے:
\begin{align*}
h(\vec x, y) & = \begin{cases}
  g(\vec x, y, h(\vec x, k(\vec x, y))) & \text{جے $k(\vec x, y) < y$} \\
  f(\vec x) & \text{نہیں تاں}
\end{cases}
\end{align*}
دوجے لفظاں وچ، $h$ دی قیمت $y$ اُتے
\emph{کسے وی} پچھلی دلیل اُتے $h$ دی قیمت ورت کے
لبھی جا سکدی اے؛ اوہ پچھلی دلیل
~$k$ دیندا اے۔

\begin{prob}
  باقی دا فنکشن $r(x,y)$ پچھلی ساریاں قیمتاں والی
  بازگشت نال متعین کرو۔ (جے $x$، $y$ قدرتی عدد
  ہون تے $y > 0$، تاں $r(x,y)$ اوہ عدد اے
  جیہڑا~$y$ توں چھوٹا اے تے
  $x = z\times y + r(x,y)$ کسے~$z$ لئی ہووے۔ گل نوں مقرر
  رکھن لئی، جے $y=0$، تاں $r(x,0) = 0$
  رکھ لَیندے آں۔)
\end{prob}

ایہہ سوچو کہ انہاں فنکشناں نوں عام بنیادی بازگشت
نال کِویں حاصل کیتا جا سکدا اے۔ بنیادی بازگشت دا
اک آخری روپ ہور لچکدار اے، کیوں جے ایس وچ
\emph{پیرامیٹر} (نال چلن والیاں قیمتاں) نوں عمل
دے دوران بدلّن دی اجازت اے:
\begin{align*}
h(\vec x, 0) & = f(\vec x) \\
h(\vec x, y+1) & = g(\vec x, y, h(k(\vec x), y))
\end{align*}
ایس دا عمل وی عام بنیادی بازگشت راہیں سرانجام دتا
جا سکدا اے۔ (ایہہ کرنا سوکھا نہیں۔ اشارے لئی، حساب
دے قدم ہتھ نال کھول کے ویکھو۔)

\end{document}
